\documentclass[11pt,a4paper]{article}

% Paquetes
\usepackage[utf8]{inputenc}
\usepackage[T1]{fontenc}
\usepackage[spanish]{babel}
\usepackage[margin=2.5cm]{geometry}
\usepackage{xcolor}
\usepackage{titlesec}
\usepackage{enumitem}
\usepackage{fancyhdr}
\usepackage{graphicx}
\usepackage{tcolorbox}
\usepackage{hyperref}

% Colores de marca
\definecolor{faramarillo}{HTML}{E0A200}
\definecolor{faranegro}{HTML}{0A0A0A}
\definecolor{faragris}{HTML}{555555}

% Enlaces
\hypersetup{
  colorlinks=true,
  linkcolor=faramarillo,
  urlcolor=faramarillo,
  citecolor=faramarillo
}

% Títulos
\titleformat{\section}
  {\color{faranegro}\Large\bfseries}
  {\thesection.}{0.5em}{}
\titleformat{\subsection}
  {\color{faramarillo}\large\bfseries}
  {\thesubsection}{0.5em}{}

% Encabezado y pie
\pagestyle{fancy}
\fancyhf{}
\lhead{\textcolor{faramarillo}{\textbf{FARA ELECTRÓNICA}}}
\rhead{Manual de uso}
\cfoot{\thepage}
\renewcommand{\headrulewidth}{0.4pt}

% Caja de nota
\newtcolorbox{nota}{
  colback=faramarillo!8,
  colframe=faramarillo,
  boxrule=0.5pt,
  arc=3pt,
  left=8pt,right=8pt,top=6pt,bottom=6pt
}

\begin{document}

% Portada
\begin{titlepage}
  \centering
  \vspace*{3cm}
  {\Huge\bfseries\textcolor{faranegro}{FARA ELECTRÓNICA}}\\[0.3cm]
  {\Large\textcolor{faramarillo}{\rule{6cm}{2pt}}}\\[1cm]
  {\LARGE Manual de Uso}\\[0.5cm]
  {\large Tienda en línea: cámaras de seguridad y calefacción}\\[2cm]
  {\large\textbf{Contenido}}\\[0.3cm]
  {\large Guía para el cliente \quad$\bullet$\quad Guía para el administrador}\\[3cm]
  {\large \today}
  \vfill
\end{titlepage}

\tableofcontents
\newpage

\section{Introducción}

Este documento explica cómo usar la tienda en línea de \textbf{Fara Electrónica},
tanto desde el punto de vista del \textbf{cliente} (la persona que compra) como del
\textbf{administrador} (quien gestiona los productos y las ventas).

La plataforma permite:
\begin{itemize}[leftmargin=*]
  \item Ver el catálogo de cámaras de seguridad y equipos de calefacción.
  \item Agregar productos a un carrito de compras.
  \item Pagar en línea con \textbf{Webpay} (tarjetas de crédito/débito) o contactar
        directamente al vendedor por \textbf{WhatsApp}.
  \item Consultar el historial de compras.
  \item (Administrador) Gestionar productos, precios, stock, ver órdenes y estadísticas
        de ventas y rentabilidad.
\end{itemize}

\begin{nota}
\textbf{Nota:} para pagar con Webpay es necesario iniciar sesión con una cuenta de Google.
El contacto por WhatsApp no requiere iniciar sesión.
\end{nota}

\section{Guía para el Cliente}

\subsection{Iniciar sesión}

\begin{enumerate}[leftmargin=*]
  \item En la esquina superior de la página, pulsa el botón \textbf{``Iniciar sesión''}.
  \item Se abrirá una ventana para entrar con tu cuenta de \textbf{Google}.
  \item Selecciona tu cuenta. Una vez dentro, verás tu nombre y tu foto en la barra superior.
\end{enumerate}

\begin{nota}
Iniciar sesión es obligatorio para pagar con Webpay y para ver tu historial de compras.
\end{nota}

\subsection{Explorar el catálogo}

Desde la página principal puedes entrar a las categorías \textbf{Cámaras de Seguridad}
o \textbf{Calefacción}. Dentro de cada categoría verás los productos en tarjetas, y podrás
ordenarlos por precio o por nombre con el menú \textbf{``Ordenar por''}.

Cada tarjeta de producto tiene tres botones:
\begin{itemize}[leftmargin=*]
  \item \textbf{Detalles:} abre una ventana con las fotos del producto (puedes ampliarlas
        pasando el cursor por encima), el precio y la descripción completa.
  \item \textbf{+ Carrito:} agrega el producto al carrito. Si no has iniciado sesión, se te
        pedirá hacerlo primero.
  \item \textbf{Comprar:} abre una ventana para comprar ese producto directamente, donde
        puedes elegir la cantidad y pagar con Webpay o contactar al vendedor.
\end{itemize}

\begin{nota}
Si un producto está \textbf{agotado}, aparecerá una etiqueta roja y no se podrá comprar.
\end{nota}

\subsection{Usar el carrito}

\begin{enumerate}[leftmargin=*]
  \item Pulsa el ícono del \textbf{carrito} en la barra superior (siempre visible).
  \item Verás la lista de productos agregados. Puedes aumentar o disminuir la cantidad
        con los botones \textbf{$+$} y \textbf{$-$}, o quitar un producto con la \textbf{X}.
  \item Abajo se muestra el \textbf{total} a pagar.
\end{enumerate}

\subsection{Pagar la compra}

En el carrito (o en la ventana ``Comprar'' de un producto) tienes dos opciones:

\textbf{Opción 1 -- Pagar con Webpay:}
\begin{enumerate}[leftmargin=*]
  \item Pulsa \textbf{``Pagar con Webpay''}.
  \item Serás redirigido al portal seguro de Webpay (Transbank).
  \item Ingresa los datos de tu tarjeta y confirma el pago.
  \item Al terminar, volverás a la página y verás si el pago fue exitoso o rechazado.
\end{enumerate}

\textbf{Opción 2 -- Contactar al vendedor por WhatsApp:}
\begin{enumerate}[leftmargin=*]
  \item Pulsa \textbf{``Contactar con vendedor''}.
  \item Se abrirá WhatsApp con un mensaje ya escrito que incluye los productos y el total.
  \item Envía el mensaje para coordinar el pago y la entrega directamente con el vendedor.
\end{enumerate}

\subsection{Ver el historial de compras}

Una vez iniciada la sesión, pulsa el ícono del \textbf{reloj} en la barra superior para
abrir \textbf{``Mis Compras''}. Ahí verás todas tus órdenes pagadas, con la fecha, los
productos, el estado (pagada, preparando, entregada, etc.) y el total.

\subsection{Modo claro / oscuro}

En la esquina inferior derecha hay un botón con un \textbf{sol/luna} que permite cambiar
entre el modo oscuro y el modo claro de la página, según tu preferencia.

\section{Guía para el Administrador}

\subsection{Acceder al panel}

El panel de administración está en la dirección:

\begin{center}
\texttt{https://(tu-dominio)/gestion-fara-super-user/}
\end{center}

Ingresa con tu \textbf{usuario y contraseña} de administrador. Tras iniciar sesión verás
el \textbf{Dashboard} con un resumen de ventas.

\begin{nota}
Solo los usuarios administradores pueden entrar a este panel y ver las estadísticas.
\end{nota}

\subsection{El Dashboard}

Al entrar verás tarjetas con las métricas principales:
\begin{itemize}[leftmargin=*]
  \item \textbf{Ventas totales:} suma de todas las órdenes pagadas.
  \item \textbf{Utilidad:} ganancia real (ventas menos los costos de los productos).
  \item \textbf{Margen:} porcentaje de rentabilidad promedio.
  \item \textbf{Órdenes:} cantidad de órdenes pagadas.
\end{itemize}
Más abajo se muestran los productos \textbf{más vendidos} y \textbf{menos vendidos},
y accesos directos a Productos, Agregar producto y Órdenes.

\subsection{Gestionar productos}

Entra a \textbf{Productos} desde el menú o el acceso directo. Allí puedes:
\begin{itemize}[leftmargin=*]
  \item Ver todos los productos con su precio, costo, margen y stock.
  \item Editar el \textbf{precio}, el \textbf{costo} y el \textbf{stock} directamente
        desde la lista.
  \item Pulsar un producto para editar todos sus datos o subir más fotos.
  \item Crear uno nuevo con el botón \textbf{``Agregar producto''}.
\end{itemize}

\begin{nota}
\textbf{Importante:} el \textbf{costo} es lo que a ti te cuesta el producto. Sirve para
calcular el margen de ganancia y \textbf{nunca se muestra al cliente}. Si no lo defines,
el margen aparecerá como 100\,\%.
\end{nota}

\subsection{Entender el margen y la utilidad}

El sistema calcula automáticamente:
\begin{itemize}[leftmargin=*]
  \item \textbf{Margen unitario} = precio de venta $-$ costo.
  \item \textbf{Margen \%} = (precio $-$ costo) $/$ precio $\times$ 100.
  \item \textbf{Utilidad de una orden} = total cobrado $-$ costo de los productos vendidos.
\end{itemize}
En la lista de productos, el margen se muestra con colores: verde (bueno),
amarillo (medio) y rojo (bajo).

\subsection{Gestionar el stock}

El stock indica cuántas unidades hay disponibles. Cuando llega a \textbf{0}, el producto
aparece como \textbf{``Agotado''} en la web y no se puede comprar. El stock se descuenta
automáticamente cuando un cliente paga con Webpay.

\subsection{Gestionar órdenes}

Entra a \textbf{Órdenes} para ver todas las compras. De cada orden puedes ver el cliente,
el total, la utilidad, el estado y los productos comprados. Puedes cambiar el \textbf{estado}
de una orden (pendiente, pagada, preparando, entregada, cancelada) para llevar el control
de las entregas.

\subsection{Agregar una nueva categoría}

Para agregar una categoría nueva de productos se requiere una pequeña modificación en el
código (en el archivo de modelos del backend) y volver a desplegar. Consulta con quien
mantiene el sistema para hacerlo de forma segura.

\section{Preguntas frecuentes}

\textbf{¿Es seguro pagar con Webpay?}\\
Sí. El pago se procesa en el portal oficial de Transbank; la tienda nunca ve ni guarda
los datos de tu tarjeta.

\vspace{0.3cm}
\textbf{¿Puedo comprar sin tarjeta?}\\
Sí, usando el botón ``Contactar con vendedor'' para coordinar por WhatsApp.

\vspace{0.3cm}
\textbf{Pagué pero no veo mi orden.}\\
Revisa ``Mis Compras''. Si el pago se descontó y la orden no aparece, contacta al vendedor
por WhatsApp con tu comprobante.

\vspace{0.3cm}
\textbf{(Admin) El margen me sale 100\,\%.}\\
Es porque el costo del producto está en 0. Edita el producto y define su costo real.

\vspace{1cm}
\begin{center}
\textcolor{faragris}{\small Manual de uso -- Fara Electrónica -- \today}
\end{center}

\end{document}
